% 52e Olympiade Internationale de Mathématiques — Amsterdam, Pays-Bas.
% Énoncés officiels en français (imo-official.org). Chaque problème vaut 7 points.

\begin{center}
  \emph{Premier jour --- lundi 18 juillet 2011 --- Durée : 4 heures 30}
\end{center}

\subsection*{Problème 1}
Pour tout ensemble $A = \{a_1, a_2, a_3, a_4\}$ de quatre entiers strictement positifs deux à deux distincts,
on note $s_A$ la somme $a_1 + a_2 + a_3 + a_4$ et on note $n_A$ le nombre de couples $(i, j)$, avec
$1 \leqslant i < j \leqslant 4$, tels que $a_i + a_j$ divise $s_A$.

Déterminer les ensembles $A$ pour lesquels $n_A$ est maximal.

\subsection*{Problème 2}
Soit $\mathcal{S}$ un ensemble fini de points du plan, contenant au moins deux points. On suppose que trois
points quelconques de $\mathcal{S}$ ne sont pas alignés.

On appelle \emph{moulin à vent} le processus suivant : le processus commence avec une droite $\ell$ contenant
un unique point $P$ de $\mathcal{S}$ ; la droite $\ell$ tourne, dans le sens des aiguilles d'une montre, autour
du point $P$, appelé \emph{pivot}, jusqu'à ce qu'elle rencontre pour la première fois un autre point de
$\mathcal{S}$ ; ce point, $Q$, devient le nouveau pivot ; la droite continue alors sa rotation dans le sens des
aiguilles d'une montre autour de $Q$, jusqu'à rencontrer un nouveau point de $\mathcal{S}$ ; ce processus
continue indéfiniment.

Montrer qu'on peut choisir un point $P$ de $\mathcal{S}$ et une droite $\ell$ contenant $P$, de façon que le
moulin à vent commençant par $\ell$ utilise chaque point de $\mathcal{S}$ comme pivot une infinité de fois.

\subsection*{Problème 3}
On désigne par $\mathbb{R}$ l'ensemble des nombres réels.

Soit $f$ une fonction de $\mathbb{R}$ dans $\mathbb{R}$ telle que, pour tous réels $x$, $y$,
\[
f(x + y) \leqslant yf(x) + f(f(x)).
\]
Montrer que $f(x) = 0$ pour tout réel $x \leqslant 0$.

\bigskip
\begin{center}
  \emph{Second jour --- mardi 19 juillet 2011 --- Durée : 4 heures 30}
\end{center}

\subsection*{Problème 4}
Soit $n$ un entier strictement positif. On dispose d'une balance à deux plateaux et de $n$ poids, de masses
respectives $2^0, 2^1, \ldots, 2^{n-1}$.

On doit placer, l'un après l'autre, chacun des $n$ poids sur la balance de telle sorte que le plateau de
droite ne soit jamais plus lourd que le plateau de gauche ; dans ce but, à chaque étape, on doit choisir un
poids qui n'est pas déjà sur la balance et le placer soit sur le plateau de gauche, soit sur le plateau de
droite ; on continue ainsi jusqu'à ce que tous les poids soient placés.

Déterminer le nombre de façons de procéder.

\subsection*{Problème 5}
On note $\mathbb{Z}$ l'ensemble des entiers et $\mathbb{N}^*$ l'ensemble des entiers strictement positifs.

Soit $f$ une fonction de $\mathbb{Z}$ dans $\mathbb{N}^*$. On suppose que, quels que soient les entiers $m$,
$n$, la différence $f(m) - f(n)$ est divisible par $f(m - n)$.

Quels que soient les entiers $m$, $n$ vérifiant $f(m) \leqslant f(n)$, montrer que $f(n)$ est divisible par
$f(m)$.

\subsection*{Problème 6}
Soit $ABC$ un triangle dont les angles sont aigus et soit $\Gamma$ son cercle circonscrit. Soit $\ell$ une
droite tangente à $\Gamma$. Soit $\ell_a$, $\ell_b$, $\ell_c$ les droites symétriques de $\ell$ par rapport
respectivement aux droites $(BC)$, $(CA)$, $(AB)$.

Montrer que le cercle circonscrit au triangle déterminé par les droites $\ell_a$, $\ell_b$, $\ell_c$ est
tangent à $\Gamma$.
